%! Justifying a Claim About the Difference of Two Means Based on a CI — Unit 7, Lesson 7.7 — Guided Notes (Answer Key)
\documentclass[10pt]{article}
\usepackage{apstats-article}
\usepackage{apstats-key}   % pulls in -boxes; do NOT also load -boxes
\newcommand{\TallMath}[1]{$\displaystyle #1\rule[-1.4em]{0pt}{3.2em}$}

\begin{document}

\pageheader{Unit 7, Lesson 7.7}{Guided Notes --- Answer Key}
\namedateperiod

\vspace{0.15in}

\begin{objectivebox}
By the end of this lesson, I will be able to:
\begin{itemize}
  \item \ansline{Interpret a CI for $\mu_1 - \mu_2$ in context, naming both populations.}
  \item \ansline{Use a CI for $\mu_1 - \mu_2$ to justify a claim by checking whether 0 is inside or outside.}
  \item \ansline{Explain how increasing sample sizes decreases the width of a two-sample CI.}
\end{itemize}
\end{objectivebox}

\begin{vocabbox}
\begin{description}
  \item[\termblanklong{Interpretation of a two-sample CI}]
        In repeated random sampling with the same sample sizes, approximately \ans{$C$}\%
        of CIs constructed this way will capture \ans{the true difference $\mu_1 - \mu_2$}.
  \item[\termblanklong{Plausible values}]
        Values \ans{inside} the CI are consistent with the data; values \ans{outside} the
        CI are not.
  \item[\termblanklong{Width of a CI}]
        $\text{width} = 2 \times t^* \cdot SE$;\quad $SE = \sqrt{\dfrac{s_1^2}{n_1} + \dfrac{s_2^2}{n_2}}$.
        Width \ans{decreases} as sample sizes increase.
\end{description}
\end{vocabbox}

\begin{notesbox}{Part 1: Interpreting a Two-Sample $t$-Interval (UNC-4.Z)}

\textbf{The Interpretation Must Include:}
\begin{enumerate}
  \item The \ans{confidence} level
  \item Both \ans{populations} (not just the samples)
  \item The \ans{direction/order} of subtraction ($\mu_1 - \mu_2$)
  \item The lower and upper bounds with \ans{units}
\end{enumerate}

\medskip
\textbf{Model Sentence (2009 AP FRQ 4):}

``Based on these samples, one can be 95 percent confident that the difference in the
population mean response times (northern $-$ southern) is between \ans{$-2.37$} minutes
and \ans{$0.37$} minutes.''

\medskip
\textbf{Worked Example.} A 99\% CI for $\mu_{\text{online}} - \mu_{\text{in-person}}$ exam
score is $(2.1, 11.3)$ points.

\medskip
Interpretation: ``We are \ans{99}\% confident that the true difference in \ans{population mean exam scores}
(\ans{online} $-$ \ans{in-person}) is between \ans{2.1} and \ans{11.3} points.''

\end{notesbox}

\begin{notesbox}{Part 2: Justifying a Claim (UNC-4.AA)}

\renewcommand{\arraystretch}{1.7}
\begin{tabularx}{\linewidth}{@{} p{0.35\linewidth} X @{}}
  \textbf{Where is 0?} & \textbf{Conclusion} \\ \hline
  0 is \emph{inside} the interval & \ansline{Insufficient evidence of a difference; a difference of 0 is plausible.} \\
  0 is \emph{above} the interval (entire interval $> 0$) & \ansline{Evidence that $\mu_1 > \mu_2$; 0 is not a plausible difference.} \\
  0 is \emph{below} the interval (entire interval $< 0$) & \ansline{Evidence that $\mu_1 < \mu_2$; 0 is not a plausible difference.} \\
\end{tabularx}

\medskip
\textbf{Example.} The 99\% CI above is $(2.1, 11.3)$. Is 0 inside or outside?

\quad 0 is \ans{below (outside)} the interval.

\quad Conclusion: There \ans{is} convincing evidence that online students score
higher than in-person students on average.

\end{notesbox}

\begin{notesbox}{Part 3: Effect of Sample Size on Width (UNC-4.AB)}

$SE = \sqrt{\dfrac{s_1^2}{n_1} + \dfrac{s_2^2}{n_2}}$. \quad As $n_1$ and $n_2$ \ans{increase},
$SE$ \ans{decreases}, so the interval becomes \ans{narrower}.

\medskip
\textbf{Example.} Suppose $s_1 = s_2 = 10$ and we compare two sample-size scenarios:

\renewcommand{\arraystretch}{1.8}
\begin{tabularx}{\linewidth}{@{} p{0.25\linewidth} X X X @{}}
  \textbf{Scenario} & $n_1 = n_2$ & $SE$ & Approx.\ margin of error \\ \hline
  A & 20 & $\sqrt{\dfrac{100}{20}+\dfrac{100}{20}} = \ans{3.16}$ & $\approx$ \ans{$2 \times 3.16 = 6.32$} \\
  B & 80 & $\sqrt{\dfrac{100}{80}+\dfrac{100}{80}} = \ans{1.58}$ & $\approx$ \ans{$2 \times 1.58 = 3.16$} \\
\end{tabularx}

\medskip
Quadrupling each sample size \ans{halves} the margin of error by a factor of \ans{2 ($\approx \sqrt{4}$)}.

\end{notesbox}

\begin{practicebox}
A 95\% CI for $\mu_{\text{brand A}} - \mu_{\text{brand B}}$ battery life is $(-3.2,\; 0.4)$ hours.

\begin{itemize}
  \item Interpret the interval: \ansline{We are 95\% confident that the true difference in mean battery life (Brand A $-$ Brand B) is between $-3.2$ and $0.4$ hours.}
  \item Does the interval provide convincing evidence that Brand A lasts longer? Explain.
        \ansline{No. 0 is inside the interval, so a difference of 0 is plausible; insufficient evidence that Brand A lasts longer.}
  \item If both sample sizes were quadrupled (all else equal), would the new CI be
        narrower, wider, or the same? Explain. \ansline{Narrower. Larger $n$ reduces $SE$, which reduces the margin of error.}
\end{itemize}
\end{practicebox}

\end{document}
